% Options for packages loaded elsewhere
\PassOptionsToPackage{unicode}{hyperref}
\PassOptionsToPackage{hyphens}{url}
%
\documentclass[
]{article}
\usepackage{amsmath,amssymb}
\usepackage{iftex}
\ifPDFTeX
  \usepackage[T1]{fontenc}
  \usepackage[utf8]{inputenc}
  \usepackage{textcomp} % provide euro and other symbols
\else % if luatex or xetex
  \usepackage{unicode-math} % this also loads fontspec
  \defaultfontfeatures{Scale=MatchLowercase}
  \defaultfontfeatures[\rmfamily]{Ligatures=TeX,Scale=1}
\fi
\usepackage{lmodern}
\ifPDFTeX\else
  % xetex/luatex font selection
\fi
% Use upquote if available, for straight quotes in verbatim environments
\IfFileExists{upquote.sty}{\usepackage{upquote}}{}
\IfFileExists{microtype.sty}{% use microtype if available
  \usepackage[]{microtype}
  \UseMicrotypeSet[protrusion]{basicmath} % disable protrusion for tt fonts
}{}
\makeatletter
\@ifundefined{KOMAClassName}{% if non-KOMA class
  \IfFileExists{parskip.sty}{%
    \usepackage{parskip}
  }{% else
    \setlength{\parindent}{0pt}
    \setlength{\parskip}{6pt plus 2pt minus 1pt}}
}{% if KOMA class
  \KOMAoptions{parskip=half}}
\makeatother
\usepackage{xcolor}
\usepackage{longtable,booktabs,array}
\usepackage{calc} % for calculating minipage widths
% Correct order of tables after \paragraph or \subparagraph
\usepackage{etoolbox}
\makeatletter
\patchcmd\longtable{\par}{\if@noskipsec\mbox{}\fi\par}{}{}
\makeatother
% Allow footnotes in longtable head/foot
\IfFileExists{footnotehyper.sty}{\usepackage{footnotehyper}}{\usepackage{footnote}}
\makesavenoteenv{longtable}
\setlength{\emergencystretch}{3em} % prevent overfull lines
\providecommand{\tightlist}{%
  \setlength{\itemsep}{0pt}\setlength{\parskip}{0pt}}
\setcounter{secnumdepth}{-\maxdimen} % remove section numbering
% unicode-fix.tex — render raw-Unicode symbols that the Latin Modern *text* font
% lacks, using Computer Modern math/text equivalents. This keeps the established
% Computer Modern look while filling the glyphs xelatex would otherwise drop.
% Included via: pandoc ... --include-in-header=unicode-fix.tex
\usepackage{newunicodechar}
\usepackage{pifont}
\usepackage{graphicx} % unicode-math's smallsetminus uses rotatebox
% Use the Computer Modern-compatible set-difference construction; the large
% Unicode glyph selected by TeX Live 2022 is absent from Latin Modern Math.
\AtBeginDocument{\let\setminus\smallsetminus}
% --- Greek letters -> math (Latin Modern Math; consistent with CM body text) ---
\newunicodechar{α}{\ensuremath{\alpha}}
\newunicodechar{β}{\ensuremath{\beta}}
\newunicodechar{γ}{\ensuremath{\gamma}}
\newunicodechar{δ}{\ensuremath{\delta}}
\newunicodechar{ε}{\ensuremath{\varepsilon}}
\newunicodechar{θ}{\ensuremath{\theta}}
\newunicodechar{μ}{\ensuremath{\mu}}
\newunicodechar{ν}{\ensuremath{\nu}}
\newunicodechar{ρ}{\ensuremath{\rho}}
\newunicodechar{σ}{\ensuremath{\sigma}}
\newunicodechar{τ}{\ensuremath{\tau}}
\newunicodechar{φ}{\ensuremath{\varphi}}
\newunicodechar{χ}{\ensuremath{\chi}}
\newunicodechar{ψ}{\ensuremath{\psi}}
\newunicodechar{η}{\ensuremath{\eta}}
\newunicodechar{κ}{\ensuremath{\kappa}}
\newunicodechar{λ}{\ensuremath{\lambda}}
\newunicodechar{π}{\ensuremath{\pi}}
\newunicodechar{ω}{\ensuremath{\omega}}
\newunicodechar{Ḣ}{\ensuremath{\dot H}}
\newunicodechar{≪}{\ensuremath{\ll}}
\newunicodechar{≫}{\ensuremath{\gg}}
\newunicodechar{⟺}{\ensuremath{\Longleftrightarrow}}
\newunicodechar{⟹}{\ensuremath{\Longrightarrow}}
\newunicodechar{⇔}{\ensuremath{\Leftrightarrow}}
\newunicodechar{⟨}{\ensuremath{\langle}}
\newunicodechar{⟩}{\ensuremath{\rangle}}
\newunicodechar{⁰}{\textsuperscript{0}}
\newunicodechar{¹}{\textsuperscript{1}}
\newunicodechar{²}{\textsuperscript{2}}
\newunicodechar{³}{\textsuperscript{3}}
\newunicodechar{⁴}{\textsuperscript{4}}
\newunicodechar{⁵}{\textsuperscript{5}}
\newunicodechar{⁶}{\textsuperscript{6}}
\newunicodechar{⁷}{\textsuperscript{7}}
\newunicodechar{⁸}{\textsuperscript{8}}
\newunicodechar{⁹}{\textsuperscript{9}}
% --- Relations / operators -> math ---
\newunicodechar{↔}{\ensuremath{\leftrightarrow}}
\newunicodechar{⇐}{\ensuremath{\Leftarrow}}
\newunicodechar{⇒}{\ensuremath{\Rightarrow}}
\newunicodechar{∂}{\ensuremath{\partial}}
\newunicodechar{∅}{\ensuremath{\varnothing}}
\newunicodechar{∈}{\ensuremath{\in}}
\newunicodechar{∼}{\ensuremath{\sim}}
\newunicodechar{≈}{\ensuremath{\approx}}
\newunicodechar{≠}{\ensuremath{\neq}}
\newunicodechar{≤}{\ensuremath{\leq}}
\newunicodechar{≥}{\ensuremath{\geq}}
\newunicodechar{⊂}{\ensuremath{\subset}}
\newunicodechar{⊃}{\ensuremath{\supset}}
% --- Special ---
\newunicodechar{ℏ}{\ensuremath{\hbar}}
\newunicodechar{ℤ}{\ensuremath{\mathbb{Z}}}
% --- Superscript modifier letters (footnote-style markers) -> text superscript ---
\newunicodechar{ʰ}{\textsuperscript{h}}
\newunicodechar{ʲ}{\textsuperscript{j}}
\newunicodechar{ᵃ}{\textsuperscript{a}}
\newunicodechar{ᵇ}{\textsuperscript{b}}
\newunicodechar{ᵈ}{\textsuperscript{d}}
\newunicodechar{ᵉ}{\textsuperscript{e}}
\newunicodechar{ᵍ}{\textsuperscript{g}}
\newunicodechar{ᶠ}{\textsuperscript{f}}
\newunicodechar{⁻}{\textsuperscript{\ensuremath{-}}}
% --- Subscript digits -> text subscript ---
\newunicodechar{₀}{\textsubscript{0}}
\newunicodechar{₁}{\textsubscript{1}}
\newunicodechar{₂}{\textsubscript{2}}
\newunicodechar{₃}{\textsubscript{3}}
\newunicodechar{₄}{\textsubscript{4}}
\newunicodechar{₈}{\textsubscript{8}}
% --- Checkmarks ---
\newunicodechar{✓}{\ding{51}}
\newunicodechar{✗}{\ding{55}}

% --- added 2026-08-02 after an independent glyph audit ---
\newunicodechar{ℓ}{\ensuremath{\ell}}
\newunicodechar{⊕}{\ensuremath{\oplus}}
\newunicodechar{⊗}{\ensuremath{\otimes}}
\newunicodechar{∠}{\ensuremath{\angle}}
\newunicodechar{≡}{\ensuremath{\equiv}}
\newunicodechar{∝}{\ensuremath{\propto}}
\newunicodechar{ℝ}{\ensuremath{\mathbb{R}}}
\newunicodechar{₅}{\ensuremath{_5}}
\newunicodechar{₋}{\ensuremath{_-}}
\newunicodechar{⁺}{\ensuremath{^+}}
\newunicodechar{⁽}{\ensuremath{^(}}
\newunicodechar{⁾}{\ensuremath{^)}}
% U+239B..U+23A0: multi-line parenthesis assembly pieces emitted by some
% converters. There is no sane inline equivalent; map to ordinary delimiters so
% the notation survives rather than vanishing silently.
\newunicodechar{⎛}{\ensuremath{\big(}}
\newunicodechar{⎜}{\ensuremath{\big\vert}}
\newunicodechar{⎝}{\ensuremath{\big(}}
\newunicodechar{⎞}{\ensuremath{\big)}}
\newunicodechar{⎟}{\ensuremath{\big\vert}}
\newunicodechar{⎠}{\ensuremath{\big)}}
\ifLuaTeX
  \usepackage{selnolig}  % disable illegal ligatures
\fi
\IfFileExists{bookmark.sty}{\usepackage{bookmark}}{\usepackage{hyperref}}
\IfFileExists{xurl.sty}{\usepackage{xurl}}{} % add URL line breaks if available
\urlstyle{same}
\hypersetup{
  hidelinks,
  pdfcreator={LaTeX via pandoc}}

\author{}
\date{}

\begin{document}

\hypertarget{non-markovian-dynamics-in-computational-biology}{%
\section{Non-Markovian Dynamics in Computational
Biology}\label{non-markovian-dynamics-in-computational-biology}}

\hypertarget{hidden-sector-memory-and-the-methodological-limits-of-single-cell-analysis}{%
\subsubsection{Hidden-Sector Memory and the Methodological Limits of
Single-Cell
Analysis}\label{hidden-sector-memory-and-the-methodological-limits-of-single-cell-analysis}}

\textbf{Author:} Alex Maybaum \textbf{Date:} May 2026
\textbf{Classification:} Computational Biology / Bioinformatics

\begin{center}\rule{0.5\linewidth}{0.5pt}\end{center}

\hypertarget{abstract}{%
\subsection{Abstract}\label{abstract}}

Single-cell bioinformatics benchmarks report setting-specific
limitations in RNA velocity, gene regulatory network inference and
perturbation prediction. Multi-modal studies and methylation-based tumor
classifiers provide distinct empirical comparisons.

We apply the finite-horizon readback diagnostic to candidate cellular
observables. A history effect at a fixed measured present gives
accessible conditional memory. Coupling, persistence and capacity
motivate possible record carriers. The cellular mechanisms and relevant
observables require further evidence.

The applications cover trajectory inference, regulatory networks,
perturbation prediction, multi-omics integration and methylation-based
classification. The benchmarks specify their data, tasks and baselines.
For instance, Heidelberg v12.8 achieved 95\% subclass accuracy in nested
cross-validation on its reference data. The mechanism and transfer to
other assays or cohorts are separate questions.

The information constraint is conditional on the target and data. A
target that differs between two completions sharing the law of the
available observations cannot be uniquely identified from that law. We
discuss where this condition could apply and how new measurements could
make an otherwise unresolved target distinguishable.

\begin{center}\rule{0.5\linewidth}{0.5pt}\end{center}

\begin{quote}
\textbf{Status --- conditional information result and conjectural
applications.} The benchmarks below are prior observations. Which
cellular targets differ across completions, which extra measurements
resolve them and which underlying mechanisms account for the comparisons
remain separate empirical questions.
\end{quote}

\hypertarget{introduction}{%
\subsection{1. Introduction}\label{introduction}}

\hypertarget{the-bioinformatics-fields-persistent-puzzles}{%
\subsubsection{1.1 The bioinformatics field's persistent
puzzles}\label{the-bioinformatics-fields-persistent-puzzles}}

Single-cell transcriptomics has revolutionized biology over the past
decade, producing datasets of millions of cells with detailed gene
expression profiles. The computational methodology has evolved rapidly:
dimensionality reduction (UMAP, t-SNE), clustering (Leiden, Louvain),
trajectory inference (Monocle, Slingshot, scVelo), gene regulatory
network inference (SCENIC, GENIE3), perturbation modeling (CPA, GEARS,
scGPT), and multi-omics integration (MOFA, Seurat, LIGER). Each new
method generation has promised improved accuracy and broader
applicability, and the field's expected trajectory has been continuous
improvement.

A different pattern has emerged. Recent benchmarking studies have
documented systematic failures across multiple method classes:

\begin{itemize}
\item
  \textbf{RNA velocity.} Gorin et al.~{[}3{]} report failures including
  trajectory reversal on particular datasets.
\item
  \textbf{Gene regulatory network inference.} Published GRN evaluations
  report limited prediction accuracy or performance close to correlation
  baselines in some settings {[}4, 5{]}.
\item
  \textbf{Deep-learning models.} A perturbation-prediction benchmark
  found no consistent improvement of its deep-learning methods over the
  simple baselines {[}6{]}. A different zero-shot study {[}7{]}
  evaluated cell clustering and other tasks.
\item
  \textbf{Multi-omics integration} methods (MOFA, MOFA+) treat different
  molecular modalities as parallel measurements of the same cell state,
  despite the fact that ``the timescale of (phospho)proteome changes
  falls into a minute regime, whereas most significant transcriptome
  changes are observed after hours'' {[}8{]}.
\end{itemize}

Each puzzle has been addressed within the field as a methodological
problem to be solved with better algorithms. The proposed solutions
follow a common pattern: incorporate additional data modalities
(chromatin accessibility, ATAC-seq, methylation), incorporate prior
knowledge (knowledge graphs, motif libraries), or move to more complex
architectures (transformers, foundation models, graph neural networks).
Some of these solutions help; others do not. The field lacks a
structural framework for predicting \emph{which} methodological
improvements will help and \emph{why}.

\hypertarget{the-oi-framework-and-its-c1c4-conditions}{%
\subsubsection{1.2 The OI framework and its C1--C4
conditions}\label{the-oi-framework-and-its-c1c4-conditions}}

The finite-horizon readback lemma gives conditional memory when
histories ending at the same present state induce different next-step
laws. The potential contributions of coupling, persistence and capacity
can be assessed for the declared partition. Those features alone allow
zero accessible memory.

The Medicine paper {[}2{]} considers chromatin, histone modifications
and other slow variables as candidate biological record carriers.
Whether measured gene expression reads back a particular retained record
is a separate test.

An identification limit arises where a particular cellular target
differs between completions sharing the entire law of the transcriptome
observations available to the method. Whether that condition holds
depends on the target, assay and interventions.

\hypertarget{thesis-and-structure}{%
\subsubsection{1.3 Thesis and structure}\label{thesis-and-structure}}

This paper applies a conditional identification bound to targets in five
domains of single-cell bioinformatics. Each application specifies the
target and available observations, considers possible identification
gaps and discusses the reported task-specific comparisons.

§2 discusses RNA and multi-omic trajectory inference. §3 considers
benchmarks for gene regulatory network inference. §4 compares
perturbation predictors on the cited tasks. §5 discusses differences in
response times across modalities. §6 reports a cancer methylome
classification example. §7 proposes variant-pathogenicity comparisons
for reader, writer and eraser genes. §8 considers bioinformatic readouts
for the disorders discussed in {[}Medicine §10.7{]}. §9 synthesizes the
cross-domain findings. §10 catalogs open questions and limits.

The bound applies when candidate completions agree on the available
observation law and disagree on the target. Whether this explains a
benchmark result requires additional evidence about the target,
observations and competing explanations for the performance difference.

\begin{center}\rule{0.5\linewidth}{0.5pt}\end{center}

\hypertarget{trajectory-inference}{%
\subsection{2. Trajectory Inference}\label{trajectory-inference}}

\hypertarget{the-rna-velocity-paradigm-and-its-assumptions}{%
\subsubsection{2.1 The RNA velocity paradigm and its
assumptions}\label{the-rna-velocity-paradigm-and-its-assumptions}}

RNA velocity, introduced by La Manno et al.~(2018) {[}9{]} and
operationalized in the velocyto and scVelo tools, infers the direction
of cellular state change from the ratio of unspliced to spliced mRNA in
single-cell RNA sequencing data. The method's appeal is that it extracts
dynamic information from static snapshots: by measuring the relative
abundance of pre-mRNA (unspliced, recently transcribed) and mature mRNA
(spliced, accumulated), it estimates whether a gene is in an induction
phase (more unspliced relative to expected steady state) or repression
phase (more spliced relative to expected steady state). Aggregating
across genes produces a velocity vector for each cell, which can be
visualized as a flow field over the dimensionally-reduced state space.

RNA velocity estimates local changes using observed spliced and
unspliced RNA under a kinetic model. The original velocyto framework and
the scVelo dynamical extension use different assumptions about steady
state and rates {[}9, 10{]}. Model adequacy depends on the tissue,
sampling and kinetic regime.

\hypertarget{documented-failure-modes}{%
\subsubsection{2.2 Documented failure
modes}\label{documented-failure-modes}}

Gorin et al.~(2022) {[}3{]} discuss examples where RNA velocity methods
fail to recover direction or reverse a known trajectory.

Heterogeneous kinetics and violations of a model's assumptions are
possible explanations for errors {[}11{]}. The relevant kinetic and
sampling conditions vary by experiment.

Flexible kinetic models and latent state models address some failures
under their own assumptions. The results depend on the available data
and benchmark.

\hypertarget{the-frameworks-structural-explanation}{%
\subsubsection{2.3 The framework's structural
explanation}\label{the-frameworks-structural-explanation}}

Chromatin state is a candidate variable in a kinetic model for RNA
trajectories. It can affect accessibility and transcription rates while
remaining latent in an RNA-only assay.

To establish accessible readback, compare later transcriptome laws for
histories with the same measured present. An association between
chromatin state and RNA measurements can motivate this test.

The sign of a velocity error depends on the model, omitted variables and
sampling design. A chromatin transition is one possible contributor.

The reported direction errors motivate testing whether adding chromatin
measurements resolves specific cases. Benchmark comparisons alone do not
isolate this cause.

\hypertarget{multi-omic-velocity-as-empirical-convergence}{%
\subsubsection{2.4 Multi-omic velocity as empirical
convergence}\label{multi-omic-velocity-as-empirical-convergence}}

MultiVelo {[}17{]} adds chromatin accessibility to RNA velocity and
reports results on the studied development and differentiation datasets.

mmVelo {[}18{]} considers single-peak representation, and MoFlow
{[}19{]} develops a multimodal velocity method. Their evaluations cover
particular data and tasks.

These studies motivate comparisons at matched targets and samples. Extra
chromatin measurements can help if they provide task-relevant
information unavailable in the RNA readout.

\hypertarget{practical-consequences}{%
\subsubsection{2.5 Practical
consequences}\label{practical-consequences}}

The framework provides three practical predictions for trajectory
inference methodology:

(P1) Test whether direction or fate is unidentifiable under the RNA-only
observation design in a setting with independently characterized
chromatin dynamics.

(P2) Compare paired RNA and chromatin data with RNA-only data at matched
benchmark conditions. A gain is plausible if the added data resolve
target differences left undetermined by RNA.

(P3) Test which measured modality adds predictive information on the
timescale of the specified trajectory.

A controlled benchmark can compare modalities and estimator classes
across prespecified trajectories.

\begin{center}\rule{0.5\linewidth}{0.5pt}\end{center}

\hypertarget{gene-regulatory-network-inference}{%
\subsection{3. Gene Regulatory Network
Inference}\label{gene-regulatory-network-inference}}

\hypertarget{the-grn-inference-paradigm-and-the-persistent-puzzle}{%
\subsubsection{3.1 The GRN inference paradigm and the ``persistent
puzzle''}\label{the-grn-inference-paradigm-and-the-persistent-puzzle}}

Gene regulatory networks (GRNs) describe the directed influence of
transcription factors on target genes. Inferring GRNs from single-cell
expression data is a foundational task in computational biology, with
applications ranging from drug-target identification to mechanistic
understanding of differentiation programs. The methodology has produced
a series of method generations: correlation-based methods (Pearson,
mutual information), ensemble methods (GENIE3, GRNBoost2),
expression-cell-type specific methods (SCENIC), causal methods (PC
algorithm, NOTEARS, DeepSEM), and most recently deep-learning foundation
models (scGPT, Geneformer) repurposed for GRN inference.

Existing comparisons report low GRN inference accuracy in some settings
{[}20{]}. Fernandez-de-Retana et al.~{[}5{]} compare causal and
correlational GRN methods under controlled benchmark conditions.

Multiple comprehensive benchmarking studies have confirmed this pattern.
McCalla et al.~{[}21{]} evaluated 13 GRN inference methods on 7
published datasets using diverse gold standards; the methods rank
similarly across gold standards, with no method showing decisive
advantage. Nguyen et al.~{[}22{]} benchmarked 15 methods on 139
simulated datasets and found extensive failure modes related to dropout
rate and sparsity. The conclusion across the literature is consistent:
GRN inference from single-cell data is hard, and current methods perform
near baseline.

\hypertarget{a-conditional-limit-on-causal-recovery}{%
\subsubsection{3.2 A conditional limit on causal
recovery}\label{a-conditional-limit-on-causal-recovery}}

Learning a causal graph from observations requires assumptions about
variables, interventions or mechanisms. Different causal structures can
represent the same observational distribution.

Chromatin state, protein concentrations and modifications are candidates
for omitted variables. Which causal target differs between
observationally equivalent completions remains open for each assay
design.

Co-expression estimates association in the observed data. Whether that
association is useful for the intended regulatory target depends on the
target and reference standard.

RNA data may leave particular causal relationships unidentified. This
does not order the predictive performance of specific correlation or
causal methods.

\hypertarget{multi-modal-methods-as-the-predicted-improvement}{%
\subsubsection{3.3 Multi-modal methods as the predicted
improvement}\label{multi-modal-methods-as-the-predicted-improvement}}

SCENIC+ {[}23{]} incorporates chromatin accessibility and motif
information to infer candidate regulatory relationships. Its reported
comparisons concern particular data and evaluation targets.

LINGER {[}24{]} uses external atlas-scale data and motif priors. The
study reports its benchmark gains under those inputs and task
definitions.

Additional observations or relevant priors could constrain the set of
compatible regulatory mechanisms. The benchmark results leave the
responsible information source and general performance ordering open.

\hypertarget{information-theoretic-ceiling}{%
\subsubsection{3.4 Information-theoretic
ceiling}\label{information-theoretic-ceiling}}

Beyond a relative improvement from multi-modal data, the framework gives
an information-theoretic ceiling. A visible law with accessible
conditional memory carries that memory in a hidden sector (the
hidden-memory theorem), and every faithful completion of the law shares
one canonical predictive quotient while the hidden fibers above it may
differ (the canonical-predictive-quotient theorem, with certified fiber
freedom) {[}Main §3.4{]}. Hence, where the causal GRN differs between
completions that share the scRNA-seq law, no method using only that law
can recover it, regardless of model complexity or training data scale.
Whether the causal structure of a given GRN lies in those fibers is
open.

scGPT {[}25{]} and Geneformer {[}26{]} offer ways to model transcriptome
associations. Kendiukhov {[}27{]} examines their attention patterns,
co-expression and prediction on the specified datasets and tasks. These
comparisons do not establish which causal relationships are identifiable
from the inputs.

Where the GRN differs across completions sharing the same available law,
scaling an RNA-only estimator cannot identify the differing part.
Evidence that the condition holds for a given target requires specifying
the observations and ruling out identification under those data.

\hypertarget{practical-consequences-1}{%
\subsubsection{3.5 Practical
consequences}\label{practical-consequences-1}}

Three framework-grounded predictions for GRN inference methodology:

(P4) Test which regulatory targets vary among completions of the same
RNA law. The bound alone does not locate existing estimators relative to
the best achievable risk.

(P5) Additional measurements can help where they resolve differences in
the specific regulatory target. Compare modalities and baselines at
fixed evaluation design.

(P6) The same identification bound applies to models using the same RNA
observation law at any scale. Paired modalities can aid recovery if they
distinguish previously equivalent completions.

The third prediction is testable in the near term as multi-modal
foundation models become available.

\begin{center}\rule{0.5\linewidth}{0.5pt}\end{center}

\hypertarget{perturbation-prediction}{%
\subsection{4. Perturbation Prediction}\label{perturbation-prediction}}

\hypertarget{the-perturbation-prediction-paradigm}{%
\subsubsection{4.1 The perturbation prediction
paradigm}\label{the-perturbation-prediction-paradigm}}

Single-cell perturbation experiments --- CRISPR knockouts (Perturb-seq,
CROP-seq), CRISPRa/i activation/repression (Replogle et al.~2022
{[}28{]}), drug treatments (sci-Plex), and combinatorial perturbations
(Norman et al.~2019 {[}29{]}) --- produce datasets where individual
cells are labeled by the perturbation they received and the resulting
transcriptional state is measured. The computational task is to predict
the transcriptional state for \emph{unseen} perturbations: combinations
of perturbations not in the training data, perturbations applied to new
cell types, or perturbations with new drugs.

This task has attracted substantial methodological investment.
Generative models (CPA {[}30{]}, chemCPA {[}31{]}) learn compositional
representations of perturbation effects. Graph neural networks (GEARS
{[}32{]}) incorporate Gene Ontology and protein-protein interaction
graphs as priors. Optimal transport methods (CellOT {[}33{]}) match
unperturbed and perturbed distributions. Most recently, single-cell
foundation models (scGPT {[}25{]}, Geneformer {[}26{]}, scFoundation
{[}34{]}, scPRINT) have been adapted for perturbation prediction tasks,
with the expectation that pretraining on tens of millions of cells would
enable strong zero-shot and few-shot performance.

\hypertarget{the-empirical-failure-of-deep-learning-on-perturbation-prediction}{%
\subsubsection{4.2 The empirical failure of deep learning on
perturbation
prediction}\label{the-empirical-failure-of-deep-learning-on-perturbation-prediction}}

Ahlmann-Eltze et al.~{[}6{]} evaluated deep-learning methods, including
a graph-based model {[}32{]}, against mean and linear baselines on the
reported perturbation tasks. The evaluated methods did not consistently
outperform the simple baselines on those datasets and metrics.

The benchmark of Fernandez-de-Retana et al.~{[}5{]} addresses GRN
inference. Kedzierska et al.~{[}7{]} evaluate separate zero-shot tasks,
and Kendiukhov {[}27{]} examines attention and regulatory signals in
specific foundation models.

These studies leave open how much of the benchmark outcome is
attributable to measurement limits, training design, task definition or
model fit.

\hypertarget{the-frameworks-structural-prediction}{%
\subsubsection{4.3 The framework's structural
prediction}\label{the-frameworks-structural-prediction}}

Let T denote the information supplied to a predictor, including the
training and intervention information it can use, and Y the defined
target. The part of a perturbation response that differs across
completions sharing the available observation law is unidentified under
that design.

If \(\widehat{Y}\) is produced from \(T\) alone, with any extra
randomness independent of \(Y\) given \(T\), data processing gives
\(I(Y;\widehat{Y})\le I(Y;T)\). For added observations \(S\), the
additional information is \(I(Y;S\mid T)\). With finite discrete
outcomes and log loss this equals the Bayes-risk reduction
\(H(Y\mid T)-H(Y\mid T,S)\). Its value depends on the design and the
target.

For squared loss, the Bayes predictor from T is the conditional mean.
The observable conditional law can also identify variance and quantiles
of the target on its measured support. A training-set mean baseline and
an optimal conditional predictor are distinct. A benchmark tie with the
former does not establish that the methods approach the latter.

\hypertarget{knowledge-graph-methods-as-the-predicted-improvement}{%
\subsubsection{4.4 Knowledge-graph methods as the predicted
improvement}\label{knowledge-graph-methods-as-the-predicted-improvement}}

GEARS {[}32{]} uses biological graphs as priors, while PerturbNet
{[}35{]} uses its own data, training and evaluation setup. Their
reported results concern the specific benchmarks. The study in {[}6{]}
also evaluated the graph-based model.

A relevant graph can supply assumptions or additional evidence for a
particular target. Performance relative to other methods remains a
benchmark question.

\hypertarget{compositional-generalization-and-what-scaling-cannot-fix}{%
\subsubsection{4.5 Compositional generalization and what scaling cannot
fix}\label{compositional-generalization-and-what-scaling-cannot-fix}}

Compositional generalization can be unidentifiable when
single-perturbation data leave multiple possible joint responses. A
specified response model or additional measurements may constrain those
alternatives.

This is testable. Norman et al.~(2019) {[}29{]} published a
combinatorial perturbation dataset with 287 single-gene knockdowns and
131 pairs of knockdowns. The framework predicts that pairs where both
genes act on the same chromatin domain (likely to interact through
hidden-sector variables) should be harder to predict from single-gene
training data than pairs where the genes act on disjoint pathways. The
empirical pattern in compositional generalization benchmarks should
track the chromatin co-localization of the perturbed genes, not just
their expression patterns. This prediction has not been directly tested
in the existing benchmarking literature.

\hypertarget{practical-consequences-2}{%
\subsubsection{4.6 Practical
consequences}\label{practical-consequences-2}}

Four framework-grounded predictions for perturbation prediction
methodology:

(P7) For an unidentified perturbation target, RNA-only methods share the
identification limit under the declared observation law. Predictive risk
for observed outcomes depends on the loss and on how well the estimator
learns the observable conditional law.

(P8) Test whether additional chromatin information improves prediction
of specified perturbation responses on paired data and matched
baselines.

(P9) Compare relevant graph or motif priors with matched data-only
methods. The effect depends on the utility of the prior for the declared
task.

(P10) Compositional generalization (predicting double perturbations from
single-perturbation training) will be systematically harder for
perturbations co-localized at the chromatin level than for perturbations
acting on disjoint pathways. This is testable through benchmarking on
Norman et al.'s combinatorial dataset stratified by chromatin
co-localization of the perturbed gene pairs.

The fourth prediction is a proposed comparison. A positive result would
have to be distinguished from alternative explanations for the
performance difference before attributing it to chromatin
co-localization.

\begin{center}\rule{0.5\linewidth}{0.5pt}\end{center}

\hypertarget{multi-omics-integration}{%
\subsection{5. Multi-Omics Integration}\label{multi-omics-integration}}

\hypertarget{the-multi-omics-integration-paradigm}{%
\subsubsection{5.1 The multi-omics integration
paradigm}\label{the-multi-omics-integration-paradigm}}

Single-cell multi-omics technologies measure multiple molecular
modalities in the same cell: G\&T-seq (genome + transcriptome),
scNMT-seq (chromatin accessibility + transcriptome + methylation),
CITE-seq (transcriptome + surface protein), QuRIE-seq (transcriptome +
intracellular phosphoprotein), and increasingly multimodal Perturb-seq
variants. The data architecture is: \(N\) cells × \(M\) modalities ×
\(G_m\) features per modality, with the modalities measuring different
aspects of cellular state on different timescales.

The standard computational task is integration: combining information
across modalities to produce unified cell embeddings, joint
trajectories, and integrated regulatory networks. The principal methods
are Multi-Omics Factor Analysis (MOFA {[}36{]} and MOFA+ {[}37{]}),
Seurat's weighted nearest neighbor (WNN) {[}38{]}, LIGER {[}39{]}, and
totalVI {[}40{]}. These methods share a common assumption: different
modalities are parallel measurements of the same underlying cell state,
and integration combines them into a single low-dimensional latent
space.

\hypertarget{the-timescale-problem}{%
\subsubsection{5.2 The timescale problem}\label{the-timescale-problem}}

The parallel-measurements assumption fails when modalities measure
different memory layers operating on different timescales. The QuRIE-seq
paper {[}8{]} explicitly states: ``The timescale of (phospho)proteome
changes falls into a minute regime, whereas most significant
transcriptome changes are observed after hours.'' The same cells,
measured at the same instant, contain information from different points
in their dynamical history --- phosphoprotein measurements reflect
minute-scale signaling state, transcriptome measurements reflect
hour-scale gene expression state, and methylation measurements reflect
day-to-week-scale chromatin state.

The framework's perspective on this is direct: each modality measures a
different memory layer with a different \(\tau_B\). The visible-sector /
hidden-sector partition is modality-dependent. From the perspective of
phosphoprotein dynamics, the transcriptome state is part of the slow
hidden sector (chromatin sets transcriptional rates, which set protein
production rates). From the perspective of methylation dynamics, the
transcriptome is part of the fast visible sector (transcription happens
fast on the methylation timescale). The static factor-analysis
assumption --- that all modalities reflect the same cell state ---
misses this hierarchical structure.

The empirical consequences are documentable. Static integration methods
produce embeddings that are dominated by the modality with the largest
dynamic range during the measurement window --- typically transcriptome
--- and the other modalities contribute mostly as confirmation of the
transcriptome-derived structure rather than as independent information.
The integration improves accuracy modestly compared to single-modality
methods, but does not fully exploit the temporal information present in
the multi-modal data.

\hypertarget{emerging-timescale-aware-methods}{%
\subsubsection{5.3 Emerging timescale-aware
methods}\label{emerging-timescale-aware-methods}}

Recent methodological work has begun to address the timescale problem
explicitly. The QuRIE-seq paper {[}8{]} used time-series sampling
combined with MOFA+ factors to capture dynamics across modalities. Heyne
et al.~(2024) {[}41{]} introduced pseudotemporal integration of
single-cell RNA-Seq and single-cell proteomics by mass spectrometry,
building cell-state trajectories separately for each modality and
aligning them through shared markers. The approach yields
``time-dependent, nonlinear relationships between transcript and protein
levels'' that static methods cannot recover.

MultiVelo {[}17{]} models chromatin and transcription with distinct
switch times and reports benchmark performance on its studied systems.

One hypothesis is that a model of finite lags improves inference where
the lagged information is present in the data and relevant to the
target. Estimating \(\tau_B\) and the shape of the response requires a
paired time-series design.

\hypertarget{the-hierarchy-of-biological-memory-layers}{%
\subsubsection{5.4 The hierarchy of biological memory
layers}\label{the-hierarchy-of-biological-memory-layers}}

The following timescales illustrate possible differences across
molecular modalities.

\begin{longtable}[]{@{}
  >{\raggedright\arraybackslash}p{(\columnwidth - 4\tabcolsep) * \real{0.3333}}
  >{\raggedright\arraybackslash}p{(\columnwidth - 4\tabcolsep) * \real{0.3333}}
  >{\raggedright\arraybackslash}p{(\columnwidth - 4\tabcolsep) * \real{0.3333}}@{}}
\toprule\noalign{}
\begin{minipage}[b]{\linewidth}\raggedright
Modality
\end{minipage} & \begin{minipage}[b]{\linewidth}\raggedright
Typical \(\tau_B\)
\end{minipage} & \begin{minipage}[b]{\linewidth}\raggedright
Memory layer
\end{minipage} \\
\midrule\noalign{}
\endhead
\bottomrule\noalign{}
\endlastfoot
Phosphoproteins & seconds--minutes & Signaling state \\
Metabolites & seconds--minutes & Metabolic state \\
Transcriptome & minutes--hours & Gene expression state \\
Histone acetylation & minutes--hours & Active chromatin state \\
Histone methylation & hours--days & Stable chromatin state \\
Proteome & hours--days & Translation state \\
Chromatin accessibility & hours--days & Regulatory element state \\
DNA methylation & days--weeks & Stable epigenetic state \\
Genome (sequence) & \(\to \infty\) & Heritable state \\
Chromatin 3D architecture (Hi-C) & hours--generations & Genome
organization \\
\end{longtable}

Timescales and their ordering can depend on the process and tissue. A
performance comparison requires estimating the relevant dynamics on the
data in question.

\hypertarget{static-factor-analysis-under-the-framework}{%
\subsubsection{5.5 Static factor analysis under the
framework}\label{static-factor-analysis-under-the-framework}}

MOFA and MOFA+ are widely used and produce useful cell embeddings. They
do not, however, capture the temporal hierarchy across modalities. The
framework's prediction is that MOFA-style methods extract the
\emph{static covariance structure} shared across modalities --- what is
common to all measurements of a cell at one instant --- but miss the
\emph{dynamical relationships} between modalities --- how chromatin
state at one time predicts transcriptome state at a later time. The
static covariance is informative but incomplete; the dynamical
relationships contain additional information that static methods cannot
access.

This is testable. A direct comparison of MOFA+ integration with a
timescale-aware method on time-series data should show the
timescale-aware method recovering dynamical relationships that MOFA
misses. The QuRIE-seq paper {[}8{]} and subsequent time-series
multi-omics studies provide the data substrate for such comparisons; the
comparison has not been systematically performed across the
methodological literature.

\hypertarget{practical-consequences-3}{%
\subsubsection{5.6 Practical
consequences}\label{practical-consequences-3}}

Three framework-grounded predictions for multi-omics integration:

(P11) For time-indexed observations, compare static factor analysis with
models of lagged dependence on the same data and prediction targets.

(P12) Test whether accounting for measured time lags improves dynamic
predictions beyond matched baselines.

(P13) Compare modality combinations at time resolutions appropriate to
the biological process and target.

Timing and modality comparisons can assess which additional measurements
or lag models aid a given task.

\begin{center}\rule{0.5\linewidth}{0.5pt}\end{center}

\hypertarget{cancer-methylome-classification}{%
\subsection{6. Cancer Methylome
Classification}\label{cancer-methylome-classification}}

\hypertarget{the-methylation-based-tumor-classification-paradigm}{%
\subsubsection{6.1 The methylation-based tumor classification
paradigm}\label{the-methylation-based-tumor-classification-paradigm}}

Capper et al.~{[}42{]} developed their CNS tumor methylation classifier
on 2,801 reference samples and reported diagnostic changes in up to 12\%
of prospective cases. Sill et al.~{[}43{]} report 95\% subclass-level
accuracy in nested cross-validation for Heidelberg v12.8, with 7,495
methylation profiles and 184 subclasses. This percentage concerns that
reference validation design.

The classifier maps measured CpG methylation profiles to curated
reference classes. Its clinical use and performance depend on reference
coverage, assay quality and the population assessed.

\hypertarget{methylome-arrest-in-cancer}{%
\subsubsection{6.2 Methylome arrest in
cancer}\label{methylome-arrest-in-cancer}}

Capper et al.~{[}42{]} discuss the potential of cell-of-origin and
tumor-associated features to inform classification. A stable signal for
some classes does not establish an unchanged global methylome during
disease.

The utility of methylation profiles for particular diagnoses motivates
longitudinal or matched-assay comparisons of which features are
retained.

\hypertarget{the-substratum-emergent-operator-distinction-in-cancer}{%
\subsubsection{6.3 The substratum-emergent operator distinction in
cancer}\label{the-substratum-emergent-operator-distinction-in-cancer}}

One hypothesis is that some methylation features retain information
about origins that is useful for a specified classification task. This
can be assessed alongside tumor-associated changes.

The 95\% result is an estimate for subclass classification within the
v12.8 reference validation design. A comparison with other assays needs
matched data and endpoints. Classification accuracy alone does not
establish which molecular features were preserved.

The chromatin and expression distinction in {[}Medicine §9.4{]}
motivates asking which assay is informative about which target in a
specified system.

\hypertarget{unclassifiable-tumors-and-substratum-disruption}{%
\subsubsection{6.4 Unclassifiable tumors and substratum
disruption}\label{unclassifiable-tumors-and-substratum-disruption}}

Drexler et al.~{[}45{]} studied Heidelberg v11.4 and used calibrated
score thresholds of 0.84 and 0.3 to define analysis groups. Tumor
purity, reference coverage and other cohort features can affect
classification.

A proposed contributor to low scores is altered methylation maintenance
in some tumors. Other explanations include low purity, technical factors
and classes absent from the reference.

A test of this hypothesis would compare prespecified maintenance-site
patterns at matched tumor purity, coverage and other clinical
covariates.

\hypertarget{methylation-classification-beyond-cns-tumors}{%
\subsubsection{6.5 Methylation classification beyond CNS
tumors}\label{methylation-classification-beyond-cns-tumors}}

Methylation-based approaches have been evaluated in additional settings,
including sarcomas {[}46{]} and pan-cancer studies {[}47{]}. Each
reported performance result uses its own reference set, validation
design and target.

An open prediction is whether changes in chromatin maintenance explain
part of the variation in classification across tumors after controlling
for reference coverage, tissue and purity.

\hypertarget{practical-consequences-4}{%
\subsubsection{6.6 Practical
consequences}\label{practical-consequences-4}}

Three framework-grounded predictions for cancer methylome
bioinformatics:

(P14) Test whether methylation features associated with specific classes
remain informative on independent samples from other tumor settings.

(P15) Compare whether maintenance-site patterns explain some low scores
after accounting for purity, reference coverage and other factors.

(P16) Test whether chromatin-maintenance covariates improve a specified
methylation classifier in held-out validation.

The third prediction is testable through reanalysis of the existing
Heidelberg cohort with maintenance machinery integrity as a
stratification variable.

\begin{center}\rule{0.5\linewidth}{0.5pt}\end{center}

\hypertarget{variant-pathogenicity-for-readerwritereraser-genes}{%
\subsection{7. Variant Pathogenicity for Reader/Writer/Eraser
Genes}\label{variant-pathogenicity-for-readerwritereraser-genes}}

\hypertarget{variant-effect-prediction-methods}{%
\subsubsection{7.1 Variant effect prediction
methods}\label{variant-effect-prediction-methods}}

Predicting the functional impact of missense variants on protein-coding
genes is a core task in clinical genetics. Methods range from
sequence-conservation methods (SIFT {[}48{]}, PolyPhen-2 {[}49{]})
through structure-aware methods (CADD {[}50{]}, REVEL {[}51{]}) to
AI-based methods (AlphaMissense {[}52{]}). AlphaMissense, introduced by
Cheng et al.~(Science 2023) {[}52{]}, applies the AlphaFold2
architecture {[}53{]} to variant effect prediction and achieves
state-of-the-art performance across multiple benchmarks, providing
pathogenicity scores for every possible single-amino-acid substitution
in the human proteome {[}54{]}.

The methods operate on a common assumption: variant pathogenicity is
determined by the variant's impact on the protein it directly affects.
Conservation features capture evolutionary intolerance to variation at
the residue; structural features capture the variant's disruption of
folding or function; AI features capture the integration of structural
and conservation information at a higher level. The methods do not, in
general, account for the variant's \emph{functional role} in the broader
regulatory architecture --- whether the affected gene is a substrate
gene, a reader of epigenetic marks, a writer that establishes marks, or
an eraser that removes marks.

\hypertarget{the-frameworks-prediction}{%
\subsubsection{7.2 The framework's
prediction}\label{the-frameworks-prediction}}

The Medicine paper proposes comparing disorders affecting readers,
writers, erasers or remodelers with disorders affecting other molecular
functions. Restoration outcomes and underlying substrate changes need to
be assessed per disorder.

The proposed variant benchmark asks whether functional grouping predicts
differences in scores or their relationship with clinical labels.

\begin{enumerate}
\def\labelenumi{(\roman{enumi})}
\item
  Test whether a declared measure of residual machinery function tracks
  graded outcomes for variants in specified genes, with appropriate
  within-gene comparisons.
\item
  Compare the relationship between variant scores and independently
  assessed phenotype severity across predeclared groups.
\item
  Compare the discrimination and, for any score interpreted as a
  probability, the calibration of predictors across matched gene groups.
\end{enumerate}

\hypertarget{the-proposed-empirical-test}{%
\subsubsection{7.3 The proposed empirical
test}\label{the-proposed-empirical-test}}

The framework's predictions can be tested through reanalysis of existing
AlphaMissense predictions stratified by gene functional role. The
protocol:

\begin{enumerate}
\def\labelenumi{\arabic{enumi}.}
\item
  Classify genes into three categories: substrate genes (encoding
  enzymes, transporters, structural proteins, ligands, receptors),
  machinery genes (encoding epigenetic readers, writers, erasers,
  chromatin remodelers, transcription factors), and architectural genes
  (encoding scaffolds and connecting structures). Standard ontologies
  (GO Molecular Function, ChromBase, Eukaryotic Linear Motif database)
  provide the classification.
\item
  For each gene category, examine the distribution of AlphaMissense
  pathogenicity scores conditional on ClinVar pathogenicity assignment.
\item
  Compare discrimination using AUC and fixed-threshold sensitivity or
  specificity at matched conservation and other covariates. If scores
  are interpreted as probabilities, assess calibration against observed
  frequencies.
\item
  Test prediction (i): within machinery genes, the variance of
  pathogenicity scores for clinically observed variants should be higher
  than within substrate genes, reflecting the graded relationship
  between machinery function and phenotype.
\item
  Test prediction (ii): correlate variant pathogenicity score with
  phenotype severity (where available from clinical literature) and
  compare slopes across gene categories.
\end{enumerate}

These comparisons can use existing public variant data. A positive
result would establish a benchmark difference under that design, with
separate evidence needed for mechanism and clinical utility.

\hypertarget{implications-for-clinical-variant-interpretation}{%
\subsubsection{7.4 Implications for clinical variant
interpretation}\label{implications-for-clinical-variant-interpretation}}

AlphaMissense supplies computational predictions {[}52{]}. Clinical
variant classification under professional standards {[}57{]} uses
multiple evidence types. A benchmark can test whether performance
differs by gene role.

A detected difference could motivate follow-up analyses with independent
functional and clinical evidence. Computational scores alone do not
establish clinical significance.

\hypertarget{practical-consequences-5}{%
\subsubsection{7.5 Practical
consequences}\label{practical-consequences-5}}

Three framework-grounded predictions for variant effect prediction
methodology:

(P17) Test whether score discrimination or probability calibration
differs across prespecified gene groups on held-out clinical labels.

(P18) Compare specialized and general predictors on the same held-out
machinery-gene data.

(P19) Test whether including gene functional role improves predictions
on an independent benchmark. Clinical classification still requires the
relevant independent evidence.

The third prediction has direct clinical implications and is testable
through retrospective audit of diagnostic genetic testing outcomes
stratified by gene category.

\begin{center}\rule{0.5\linewidth}{0.5pt}\end{center}

\hypertarget{reader-writer-disorder-bioinformatics}{%
\subsection{8. Reader-Writer Disorder
Bioinformatics}\label{reader-writer-disorder-bioinformatics}}

\hypertarget{the-frameworks-reader-writer-disorder-class}{%
\subsubsection{8.1 The framework's reader-writer disorder
class}\label{the-frameworks-reader-writer-disorder-class}}

The proposed comparisons span distinct mechanisms: MECP2 binding to
methylated DNA, FMR1 RNA binding, UBE3A ubiquitin ligase activity, KMT2D
methyltransferase activity, KDM6A H3K27 demethylation and chromatin
remodeling in ATR-X syndrome. A shared pattern of retained and altered
molecular features is a hypothesis.

\hypertarget{bioinformatic-patterns-in-rett-syndrome}{%
\subsubsection{8.2 Bioinformatic patterns in Rett
syndrome}\label{bioinformatic-patterns-in-rett-syndrome}}

Rett syndrome provides a case for comparing affected readouts. MECP2
binds methylated DNA. The extent and interpretation of global
methylation changes in MECP2 loss-of-function remain open {[}60{]}.

Picard and Fagiolini {[}60{]} review unresolved effects of MeCP2
malfunction on DNA methylation. Qian et al.~{[}59{]} studied restoration
in a human neuronal model.

Transcriptional changes described in MECP2 models include changes
affecting long genes {[}61{]}. These results motivate paired measurement
of methylation and transcriptome readouts under the specified
conditions.

\hypertarget{generalization-to-other-reader-writer-disorders}{%
\subsubsection{8.3 Generalization to other reader-writer
disorders}\label{generalization-to-other-reader-writer-disorders}}

The proposed pattern can be assessed in the following disorders with
distinct molecular functions.

\textbf{Fragile X syndrome (FMR1).} FMR1 encodes an RNA-binding protein
that regulates translation. Compare RNA and protein readouts in a
specified model to assess the effects of altered FMR1 function.

\textbf{Angelman syndrome (UBE3A).} UBE3A encodes an E3 ubiquitin
ligase. Joint protein, ubiquitination and RNA readouts can test which
effects follow loss or restoration under a specified condition.

\textbf{Kabuki syndrome (KMT2D and KDM6A).} KMT2D is an H3K4
methyltransferase and KDM6A removes H3K27 methylation. Whether other
features are retained and which respond to restoration depend on the
disorder and model.

\textbf{ATR-X syndrome (ATRX).} ATRX remodels chromatin at sites
including repetitive DNA. The distribution of altered and retained
features is an empirical question.

Whether a profile of retained and altered features generalizes across
these disorders remains open.

\hypertarget{differential-diagnosis-through-multi-omic-bioinformatics}{%
\subsubsection{8.4 Differential diagnosis through multi-omic
bioinformatics}\label{differential-diagnosis-through-multi-omic-bioinformatics}}

A proposed differential analysis would compare paired molecular readouts
under matched sampling and clinical conditions.

\begin{enumerate}
\def\labelenumi{\arabic{enumi}.}
\item
  \textbf{DNA methylation.} Compare the measured loci and global
  profiles under a declared assay.
\item
  \textbf{Histone modifications.} Compare relevant marks at prespecified
  loci.
\item
  \textbf{Chromatin accessibility.} Test for locus-level or broader
  differences while accounting for cell composition.
\item
  \textbf{Transcriptome.} Test for reproducible expression differences
  under the specified disorder and model.
\end{enumerate}

The usefulness of a combined profile for differential diagnosis requires
validation against known diagnoses in independent data.

\hypertarget{therapeutic-implications-and-bioinformatic-monitoring}{%
\subsubsection{8.5 Therapeutic implications and bioinformatic
monitoring}\label{therapeutic-implications-and-bioinformatic-monitoring}}

Mecp2 reactivation in Rett mice produced neurological rescue (Guy et
al.~2007, {[}Medicine §10.7{]}). This motivates examining biomarkers of
restoration. Class-wide substrate preservation and wider treatment
windows remain open. Trofinetide is a separate approved Rett treatment
with unknown mechanism of therapeutic effect {[}67{]}.

\begin{enumerate}
\def\labelenumi{(\roman{enumi})}
\item
  Test how expression and methylation change after restoration in a
  specified model, with measured levels of restored function.
\item
  Test whether baseline methylation predicts response among matched
  cases alongside alternative reasons for nonresponse.
\item
  Compare candidate molecular and functional markers over follow-up to
  assess whether they predict sustained response.
\end{enumerate}

\hypertarget{practical-consequences-6}{%
\subsubsection{8.6 Practical
consequences}\label{practical-consequences-6}}

Three framework-grounded predictions for reader-writer disorder
bioinformatics:

(P20) Test whether a proposed profile of retained marks and altered
expression generalizes across named disorders.

(P21) Evaluate whether multi-omic profiles distinguish defined disorder
groups on an independent diagnostic cohort.

(P22) Test whether a specified baseline substrate measure predicts
response to restoration of function.

The second and third predictions have direct clinical applications and
motivate prospective multi-omic studies in reader-writer disorder
cohorts.

\begin{center}\rule{0.5\linewidth}{0.5pt}\end{center}

\hypertarget{synthesis}{%
\subsection{9. Synthesis}\label{synthesis}}

\hypertarget{empirical-comparisons-and-proposed-mechanisms}{%
\subsubsection{9.1 Empirical comparisons and proposed
mechanisms}\label{empirical-comparisons-and-proposed-mechanisms}}

The preceding sections have developed framework predictions across five
active methodological domains in single-cell bioinformatics. The
findings are summarized in Table 1.

\begin{longtable}[]{@{}
  >{\raggedright\arraybackslash}p{(\columnwidth - 8\tabcolsep) * \real{0.2000}}
  >{\raggedright\arraybackslash}p{(\columnwidth - 8\tabcolsep) * \real{0.2000}}
  >{\raggedright\arraybackslash}p{(\columnwidth - 8\tabcolsep) * \real{0.2000}}
  >{\raggedright\arraybackslash}p{(\columnwidth - 8\tabcolsep) * \real{0.2000}}
  >{\raggedright\arraybackslash}p{(\columnwidth - 8\tabcolsep) * \real{0.2000}}@{}}
\toprule\noalign{}
\begin{minipage}[b]{\linewidth}\raggedright
Domain
\end{minipage} & \begin{minipage}[b]{\linewidth}\raggedright
Reported comparison
\end{minipage} & \begin{minipage}[b]{\linewidth}\raggedright
Proposed factor to test
\end{minipage} & \begin{minipage}[b]{\linewidth}\raggedright
Candidate additional input
\end{minipage} & \begin{minipage}[b]{\linewidth}\raggedright
Scope of evidence
\end{minipage} \\
\midrule\noalign{}
\endhead
\bottomrule\noalign{}
\endlastfoot
Trajectory inference (§2) & Direction errors in named datasets {[}3{]} &
Omitted kinetics or chromatin state & Multi-omic velocity &
Method-specific reports {[}17--19{]} \\
GRN inference (§3) & Variable performance relative to correlation
{[}5{]} & Different GRNs consistent with RNA data & Modalities and
independent priors & Benchmarks for SCENIC+ {[}23{]} and LINGER
{[}24{]} \\
Perturbation prediction (§4) & Deep models compared with mean and linear
baselines {[}6{]} & Target differences invisible to specified inputs &
Relevant external observations or priors & Design-specific
performance \\
Multi-omics integration (§5) & Modalities can have different response
times {[}8{]} & Lag structure & Time-indexed multi-omics data &
Method-specific evaluations {[}8, 41{]} \\
Cancer methylome (§6) & 95\% subclass accuracy for Heidelberg v12.8
{[}43{]} & Which features carry diagnostic information & Additional
assays on matched cases & Nested cross-validation for the specified
classifier \\
\end{longtable}

A shared question is whether the available data distinguish the target
under study. Task-specific comparisons motivate testing that question
alongside alternative explanations.

\hypertarget{the-information-theoretic-ceiling}{%
\subsubsection{9.2 The information-theoretic
ceiling}\label{the-information-theoretic-ceiling}}

A quantity that differs between completions sharing the law of all
observations available to the method is unidentified from that law. The
target and observational design determine whether this limit applies.

The quantities proposed as subject to the ceiling include: the cell's
causal regulatory network structure (§3), the response to unseen
perturbations (§4), the temporal lag relationships between chromatin and
transcription (§5), and the substratum-level memory operator (§6). For
each, the underlying biology integrates over hidden-sector variables
before the data is collected; whether the quantity lies in the hidden
fibers is open.

Foundation-model benchmarks {[}6, 7{]} motivate testing candidate
identification gaps. Their scores also depend on estimators, training
data and evaluation design.

If the gap holds, resolving the unidentified part requires information
or independently justified assumptions that distinguish the candidate
completions.

\hypertarget{the-roadmap-forward}{%
\subsubsection{9.3 The roadmap forward}\label{the-roadmap-forward}}

The following approaches can be compared on targets with specified
observation limits.

\textbf{Pathway 1: Multi-modal data.} Chromatin accessibility (10x
Multiome, scNMT-seq), histone modifications (CITE-seq variants,
Paired-Tag), methylation (scNMT-seq) and chromatin architecture (scHi-C)
provide measurements beyond the transcriptome. The reported comparisons
of paired-data methods, including MultiVelo and the multiomic inference
method in {[}23{]}, concern specific tasks, datasets and baselines.
Benefit for another target depends on the additional information and how
the estimator uses it.

\textbf{Pathway 2: Prior biological knowledge.} Knowledge graphs (Gene
Ontology, KEGG pathways, protein-protein interaction networks), motif
libraries (JASPAR, HOCOMOCO) and curated regulatory databases (Regulome,
ENCODE TFBS) provide candidate priors from other sources. GEARS uses
such constraints. SCENIC likewise uses regulatory priors. The
atlas-scale data and motif priors in {[}24{]} provide another example.
Their reported performance is specific to the task, data and comparator.

\textbf{Pathway 3: Explicit memory state representation.} Models that
maintain explicit chromatin-state variables alongside transcriptome
variables --- separate from the latent representations of deep learning
models --- can in principle capture non-Markovian dynamics. This
approach is less developed but theoretically aligned with the
framework's prediction. Mechanistic models in the MultiVelo style
(coupled ODEs for chromatin and transcription) are early examples.

The three pathways are complementary, not competing. The empirical
pattern suggests that combining them --- multi-modal data analyzed with
knowledge-graph priors using methods that maintain explicit memory state
--- represents the most promising methodological direction.

\hypertarget{measurements-of-persistent-state-and-expression}{%
\subsubsection{9.4 Measurements of persistent state and
expression}\label{measurements-of-persistent-state-and-expression}}

Chromatin, DNA methylation and protein assays can measure features left
unobserved by RNA assays. Whether these measurements resolve a
susceptibility, causal or predictive target depends on the design.

The operator distinction motivates the comparison of molecular state
with downstream expression and function. The empirical studies test
particular readouts.

Methylation-based classification and specific restoration studies
illustrate different applications. Neither establishes a common physical
mechanism or class-wide preservation of a molecular substrate.

\begin{center}\rule{0.5\linewidth}{0.5pt}\end{center}

\hypertarget{open-questions-and-limits}{%
\subsection{10. Open Questions and
Limits}\label{open-questions-and-limits}}

\hypertarget{what-the-framework-predicts-but-has-not-been-tested}{%
\subsubsection{10.1 What the framework predicts but has not been
tested}\label{what-the-framework-predicts-but-has-not-been-tested}}

Several framework predictions in this paper are testable but not yet
empirically verified:

\begin{enumerate}
\def\labelenumi{(\roman{enumi})}
\item
  \textbf{Variant benchmark} (§7). Compare classification,
  discrimination and any probability calibration across prespecified
  gene groups.
\item
  \textbf{Compositional generalization and chromatin co-localization}
  (§4.5). The prediction that double-perturbation prediction difficulty
  correlates with the chromatin co-localization of the perturbed genes
  can be tested using existing Perturb-seq datasets stratified by
  chromatin features.
\item
  \textbf{Unclassifiable tumor substratum disruption} (§6.4). The
  prediction that unclassifiable CNS tumors have specifically altered
  methylation at maintenance sites versus content sites is testable in
  the existing Heidelberg dataset.
\item
  \textbf{Timescale-aware integration} (§5.6). Test whether measured lag
  structure improves specified tasks on controlled benchmark data.
\end{enumerate}

Each is a discrete empirical test with clear positive and negative
outcomes. The framework's bioinformatic predictions are falsifiable in a
specific operational sense.

\hypertarget{what-the-framework-does-not-predict}{%
\subsubsection{10.2 What the framework does not
predict}\label{what-the-framework-does-not-predict}}

The framework's predictions are limited to information-theoretic
constraints and structural patterns. It does not predict:

\begin{enumerate}
\def\labelenumi{(\roman{enumi})}
\item
  \textbf{Algorithm performance.} The conditional identification limit
  alone does not rank named estimators. Additional measurements can
  lower optimal risk when relevant to the target and loss.
\item
  \textbf{Effect sizes.} Benefits from a knowledge graph depend on the
  source of the prior, its relevance and the estimator.
\item
  \textbf{Failure modes outside the ceiling-bounded tasks.} The
  framework's information-theoretic ceiling applies to tasks requiring
  hidden-sector information. Other failure modes (batch effects,
  technical noise, dropouts, low cell numbers) are real and important
  but outside the framework's prediction scope.
\item
  \textbf{Specific molecular mechanisms.} The framework predicts
  structural patterns at the level of memory layers and information
  flows, not specific molecular pathways. Mechanistic biology continues
  to require domain-specific theoretical and experimental work.
\end{enumerate}

\hypertarget{open-structural-questions}{%
\subsubsection{10.3 Open structural
questions}\label{open-structural-questions}}

Three open questions remain at the structural level:

\textbf{Q1. What is the empirical \(\tau_B\) hierarchy across biological
systems?} The framework predicts that different memory layers have
different characteristic timescales, with consequences for which
modalities matter for which questions (§5.4). The actual \(\tau_B\)
values are tissue- and process-dependent, and have not been
systematically characterized. A program of measurements characterizing
\(\tau_B\) across cell types, tissues, and biological processes would
substantially refine the framework's bioinformatic predictions.

\textbf{Q2. How do different machinery genes contribute to the
trace-out?} The framework treats epigenetic readers, writers, and
erasers as components of the trace-out machinery, but the specific
contribution of each to the substratum-to-emergent translation is not
characterized in detail. Different machinery genes may produce different
patterns of substratum-emergent dissociation, with consequences for the
differential bioinformatic signatures of different reader-writer
disorders (§8.3).

\textbf{Q3. What is the relationship between the framework's
bioinformatic predictions and explicit kinetic models?} The framework's
predictions are information-theoretic and structural. Specific kinetic
models (mass-action models of gene regulation, stochastic gene
expression models, polymer models of chromatin) are mechanistic. The
relationship between the framework's structural picture and the
mechanistic models has not been worked out in detail. A formal
connection would clarify which aspects of mechanistic modeling
correspond to which framework structural features.

\hypertarget{limits}{%
\subsubsection{10.4 Limits}\label{limits}}

The framework's bioinformatic predictions have inherent limits:

\begin{enumerate}
\def\labelenumi{(\roman{enumi})}
\item
  \textbf{Scope of conditional memory.} C4 concerns a specified
  matched-present history effect. Other data limitations, such as
  unmeasured variables or poor reference coverage, can constrain a
  prediction independently of C4.
\item
  \textbf{Attainable performance.} Estimation error, model design and
  limited training data can keep a method below the performance
  achievable from its inputs.
\item
  \textbf{Empirical validation.} A comparison needs a specified target,
  data, loss and controls. Confounders include batch effects, dataset
  selection and evaluation design.
\item
  \textbf{The framework does not replace mechanistic biology.}
  Bioinformatic methods aim to extract biological insight from data, and
  biological insight requires understanding mechanisms. The framework
  provides a structural perspective on what mechanisms determine method
  performance, but does not substitute for the mechanistic understanding
  itself.
\end{enumerate}

The mathematical limit is conditional on observational equivalence and
target differences. The cellular mechanisms and relative performance
proposals remain separate tests.

\begin{center}\rule{0.5\linewidth}{0.5pt}\end{center}

\hypertarget{acknowledgements}{%
\subsection{Acknowledgements}\label{acknowledgements}}

AI assistance in preparing this work is disclosed at\\
\url{https://github.com/amaybaum/incompleteness\#ai-assistance}.

\begin{center}\rule{0.5\linewidth}{0.5pt}\end{center}

\hypertarget{references}{%
\subsection{References}\label{references}}

{[}1{]} A. Maybaum, ``The Incompleteness of Observation,'' (2026).

{[}2{]} A. Maybaum, ``Non-Markovian Dynamics in Biology and Medicine,''
companion paper (2026).

{[}3{]} G. Gorin, M. Fang, T. Chari, and L. Pachter, ``RNA velocity
unraveled,'' \emph{PLOS Computational Biology} \textbf{18}, e1010492
(2022).

{[}4{]} X. Yuan et al., ``Inferring gene regulatory networks from
single-cell multiome data using atlas-scale external data,''
\emph{Nature Biotechnology} (2024). doi:10.1038/s41587-024-02182-7.

{[}5{]} M. Fernandez-de-Retana, R. Sanchez-Corcuera, U. Zulaika, A.
Bilbao-Jayo, and A. Almeida, ``When Does Gene Regulatory Network
Inference Break? A Controlled Diagnostic Study of Causal and
Correlational Methods on Single-Cell Data,'' arXiv:2605.04930 (2026).

{[}6{]} C. Ahlmann-Eltze et al., ``Deep-learning-based gene perturbation
effect prediction does not yet outperform simple linear baselines,''
\emph{Nature Methods} (2025). doi:10.1038/s41592-025-02772-6.

{[}7{]} K. Kedzierska et al., ``Zero-shot evaluation reveals limitations
of single-cell foundation models,'' \emph{Genome Biology} 26, 101
(2025). doi:10.1186/s13059-025-03574-x.

{[}8{]} M. Reimegård et al., ``Single-cell intracellular epitope and
transcript detection reveals signal transduction dynamics'' (QuRIE-seq),
\emph{Cell Reports Methods} 1, 100165 (2021).

{[}9{]} G. La Manno et al., ``RNA velocity of single cells,''
\emph{Nature} 560, 494--498 (2018).

{[}10{]} V. Bergen et al., ``Generalizing RNA velocity to transient cell
states through dynamical modeling'' (scVelo), \emph{Nature
Biotechnology} 38, 1408--1414 (2020).

{[}11{]} M. Barile et al., ``Coordinated changes in gene expression
kinetics underlie both mouse and human erythroid maturation,''
\emph{Genome Biology} 22, 197 (2021).

{[}12{]} H. Cui et al., ``DeepVelo: Deep Learning extends RNA velocity
to multi-lineage systems with cell-specific kinetics,'' bioRxiv (2022).

{[}13{]} L. Liang et al., ``LatentVelo: structured latent gene
expression dynamics for RNA velocity estimation,'' bioRxiv (2022).

{[}14{]} A. Gayoso et al., ``Deep generative modeling of transcriptional
dynamics for RNA velocity analysis in single cells'' (veloVI),
\emph{Nature Methods} 21, 50--59 (2024).

{[}15{]} L. Liang et al., ``LatentVelo: structured latent gene
expression dynamics,'' bioRxiv (2022).

{[}16{]} A. Gayoso et al., ``veloVI: deep generative modeling of
transcriptional dynamics,'' \emph{Nature Methods} 21, 50--59 (2024).

{[}17{]} C. Li et al., ``Multi-omic single-cell velocity models
epigenome-transcriptome interactions and improves cell fate prediction''
(MultiVelo), \emph{Nature Biotechnology} 40, 1690--1698 (2022).

{[}18{]} N. Yamanaka et al., ``mmVelo: A deep generative model for
estimating cell state-dependent dynamics across multiple modalities,''
bioRxiv (2024). doi:10.1101/2024.12.11.628059.

{[}19{]} T. Tang et al., ``MoFlow: Multi-omic relay velocity modeling
uncovers dynamic chromatin-transcription regulation across cell
states,'' \emph{Nature Communications} (2025).
doi:10.1038/s41467-025-67259-6.

{[}20{]} X. Yuan et al., ``Inferring gene regulatory networks from
single-cell multiome data using atlas-scale external data,''
\emph{Nature Biotechnology} 42, 247--257 (2024).

{[}21{]} S. McCalla et al., ``Comprehensive Benchmarking of Single-Cell
Gene Regulatory Network Inference,'' \emph{Genome Biology} 24, 21
(2023).

{[}22{]} H. Nguyen et al., ``A comprehensive survey of regulatory
network inference methods using single cell RNA sequencing data,''
\emph{Briefings in Bioinformatics} 22, bbaa190 (2021).

{[}23{]} C. Bravo González-Blas et al., ``SCENIC+: single-cell multiomic
inference of enhancers and gene regulatory networks,'' \emph{Nature
Methods} 20, 1355--1367 (2023).

{[}24{]} X. Yuan et al., ``LINGER: Lifelong neural network for gene
regulation from single-cell multiome data,'' \emph{Nature Biotechnology}
42, 247--257 (2024).

{[}25{]} H. Cui et al., ``scGPT: toward building a foundation model for
single-cell multi-omics using generative AI,'' \emph{Nature Methods} 21,
1470--1480 (2024).

{[}26{]} C. V. Theodoris et al., ``Transfer learning enables predictions
in network biology'' (Geneformer), \emph{Nature} 618, 616--624 (2023).

{[}27{]} I. Kendiukhov, ``Systematic Evaluation of Single-Cell
Foundation Model Interpretability Reveals Attention Captures
Co-Expression Rather Than Unique Regulatory Signal,'' arXiv:2602.17532
(2026).

{[}28{]} J. M. Replogle et al., ``Mapping information-rich
genotype-phenotype landscapes with genome-scale Perturb-seq,''
\emph{Cell} 185, 2559--2575 (2022).

{[}29{]} T. M. Norman et al., ``Exploring genetic interaction manifolds
constructed from rich single-cell phenotypes,'' \emph{Science} 365,
786--793 (2019).

{[}30{]} M. Lotfollahi et al., ``Predicting cellular responses to
complex perturbations in high-throughput screens'' (CPA),
\emph{Molecular Systems Biology} 19, e11517 (2023).

{[}31{]} L. Hetzel et al., ``chemCPA: Predicting Cellular Responses to
Novel Drug Perturbations at a Single-Cell Resolution,'' NeurIPS (2022).

{[}32{]} Y. Roohani, K. Huang, J. Leskovec, ``Predicting transcriptional
outcomes of novel multigene perturbations with GEARS,'' \emph{Nature
Biotechnology} 42, 927--935 (2024).

{[}33{]} C. Bunne et al., ``Learning single-cell perturbation responses
using neural optimal transport,'' \emph{Nature Methods} 20, 1759--1768
(2023).

{[}34{]} M. Hao et al., ``Large-scale foundation model on single-cell
transcriptomics'' (scFoundation), \emph{Nature Methods} 21, 1481--1491
(2024).

{[}35{]} H. Yu et al., ``PerturbNet predicts single-cell responses to
unseen chemical and genetic perturbations,'' \emph{Molecular Systems
Biology} (2025).

{[}36{]} R. Argelaguet et al., ``Multi-Omics Factor Analysis (MOFA) ---
a framework for unsupervised integration of multi-omics data sets,''
\emph{Molecular Systems Biology} 14, e8124 (2018).

{[}37{]} R. Argelaguet et al., ``MOFA+: a statistical framework for
comprehensive integration of multi-modal single-cell data,''
\emph{Genome Biology} 21, 111 (2020).

{[}38{]} Y. Hao et al., ``Integrated analysis of multimodal single-cell
data'' (Seurat v4), \emph{Cell} 184, 3573--3587 (2021).

{[}39{]} J. D. Welch et al., ``Single-cell multi-omic integration
compares and contrasts features of brain cell identity'' (LIGER),
\emph{Cell} 177, 1873--1887 (2019).

{[}40{]} A. Gayoso et al., ``Joint probabilistic modeling of single-cell
multi-omic data with totalVI,'' \emph{Nature Methods} 18, 272--282
(2021).

{[}41{]} R. Heyne et al., ``Resolving Single-Cell Gene Expression by
Pseudotemporal Integration of Transcriptomic and Proteomic Datasets,''
\emph{PLOS Computational Biology} (2024).

{[}42{]} D. Capper et al., ``DNA methylation-based classification of
central nervous system tumours,'' \emph{Nature} 555, 469--474 (2018).

{[}43{]} M. Sill et al., ``Advancing CNS tumor diagnostics with expanded
DNA methylation-based classification'' (Heidelberg classifier v12.8),
\emph{Cancer Cell}. doi:10.1016/j.ccell.2025.11.002.

{[}44{]} D. N. Louis et al., ``The 2021 WHO Classification of Tumors of
the Central Nervous System: a summary,'' \emph{Neuro-Oncology} 23,
1231--1251 (2021).

{[}45{]} R. Drexler et al., ``Unclassifiable CNS tumors in DNA
methylation-based classification: clinical challenges and prognostic
impact,'' \emph{Acta Neuropathologica Communications} (2024).
doi:10.1186/s40478-024-01728-9.

{[}46{]} C. Koelsche et al., ``Sarcoma classification by DNA methylation
profiling,'' \emph{Nature Communications} 12, 498 (2021).

{[}47{]} M. Renaud et al., ``Pan-cancer DNA methylation
classification,'' \emph{Genome Medicine} 14, 80 (2022).

{[}48{]} P. C. Ng, S. Henikoff, ``SIFT: predicting amino acid changes
that affect protein function,'' \emph{Nucleic Acids Research} 31,
3812--3814 (2003).

{[}49{]} I. A. Adzhubei et al., ``A method and server for predicting
damaging missense mutations'' (PolyPhen-2), \emph{Nature Methods} 7,
248--249 (2010).

{[}50{]} M. Kircher et al., ``A general framework for estimating the
relative pathogenicity of human genetic variants'' (CADD), \emph{Nature
Genetics} 46, 310--315 (2014).

{[}51{]} N. M. Ioannidis et al., ``REVEL: An Ensemble Method for
Predicting the Pathogenicity of Rare Missense Variants,'' \emph{American
Journal of Human Genetics} 99, 877--885 (2016).

{[}52{]} J. Cheng et al., ``Accurate proteome-wide missense variant
effect prediction with AlphaMissense,'' \emph{Science} 381, eadg7492
(2023).

{[}53{]} J. Jumper et al., ``Highly accurate protein structure
prediction with AlphaFold,'' \emph{Nature} 596, 583--589 (2021).

{[}54{]} L. Geistlinger et al., ``AlphaMissenseR: an integrated
framework for investigating missense mutations in human protein-coding
genes,'' \emph{Bioinformatics} (2025).
doi:10.1093/bioinformatics/btaf093.

{[}55{]} J. Guy et al., ``Reversal of neurological defects in a mouse
model of Rett syndrome,'' \emph{Science} 315, 1143--1147 (2007).

{[}56{]} D. C. Tarquinio et al., ``The course of awake breathing
disturbances across the lifespan in Rett syndrome,'' \emph{Brain \&
Development} 40, 515--529 (2018).

{[}57{]} S. Richards et al., ``Standards and guidelines for the
interpretation of sequence variants: a joint consensus recommendation of
the American College of Medical Genetics and Genomics and the
Association for Molecular Pathology,'' \emph{Genetics in Medicine} 17,
405--423 (2015).

{[}58{]} J. Kaufmann et al., ``Expert Recommendations for Trofinetide
Use in Individuals with Rett Syndrome,'' \emph{Pediatric Neurology}
(2024).

{[}59{]} J. Qian et al., ``Multiplex epigenome editing of MECP2 to
rescue Rett syndrome neurons,'' \emph{Science Translational Medicine}
15, eadd4666 (2023).

{[}60{]} N. Picard, M. Fagiolini, ``Transcriptomic and Epigenomic
Landscape in Rett Syndrome,'' \emph{Genes} 12, 967 (2021).

{[}61{]} H. W. Gabel et al., ``Disruption of DNA-methylation-dependent
long gene repression in Rett syndrome,'' \emph{Nature} 522, 89--93
(2015).

{[}62{]} J. C. Darnell et al., ``FMRP stalls ribosomal translocation on
mRNAs linked to synaptic function and autism,'' \emph{Cell} 146,
247--261 (2011).

{[}63{]} M. R. Sawicka et al., ``FMRP has a cell-type-specific role in
CA1 pyramidal neurons to regulate autism-related transcripts and
circadian memory,'' \emph{eLife} 8, e46919 (2019).

{[}64{]} Y. Yi et al., ``An autism-linked mutation disables
phosphorylation control of UBE3A,'' \emph{Cell} 162, 795--807 (2015).

{[}65{]} R. C. Cuellar et al., ``KMT2D and chromatin in development and
disease,'' \emph{Genes \& Development} 36, 1083--1100 (2022).

{[}66{]} G. Goldberg et al., ``Distinct factors control histone variant
H3.3 localization at specific genomic regions,'' \emph{Cell} 140,
678--691 (2010).

{[}67{]} U.S. prescribing information for trofinetide,
\href{https://dailymed.nlm.nih.gov/dailymed/fda/fdaDrugXsl.cfm?setid=67e6f2d9-21f6-466f-9def-826c6a4b8257}{§12.1
Mechanism of Action}.

\end{document}
% source-sha256: 2051c86b5c61559ecea8c92e0f0089cf13571c187880db4b88381d5f9bee25f9
